% ==============================================================================
% ПРЕАМБУЛА: Учебно-методическое пособие по дисциплине
% «Физические основы нанотехнологий» (ГУАП, Кафедра №3)
% ==============================================================================

\usepackage[T2A]{fontenc}
\usepackage[utf8]{inputenc}
\usepackage[russian]{babel}

% Геометрия страницы по университетскому стандарту
\usepackage[a4paper, top=20mm, bottom=20mm, left=25mm, right=15mm, headheight=15pt]{geometry}

% Нумерация в стиле «Глава 2» (2.X)
\renewcommand{\thesection}{2.\arabic{section}}
\renewcommand{\thesubsection}{2.\arabic{section}.\arabic{subsection}}
\renewcommand{\thesubsubsection}{2.\arabic{section}.\arabic{subsection}.\arabic{subsubsection}}
\renewcommand{\theequation}{2.\arabic{equation}}
\renewcommand{\thefigure}{2.\arabic{figure}}
\renewcommand{\thetable}{2.\arabic{table}}

% Типографика и микротипографика
\usepackage{microtype}
\usepackage{indentfirst}
\setlength{\parindent}{1.25cm}
\setlength{\parskip}{2pt plus 1pt minus 1pt}
\linespread{1.15}

% Математические пакеты
\usepackage{amsmath,amssymb,mathtools,bm}

% Таблицы
\usepackage{booktabs}
\usepackage{tabularx}
\usepackage{array}
\usepackage{multirow}

% Графика и подрисунки
\usepackage{graphicx}
\usepackage{caption}
\usepackage{subcaption}
\captionsetup{font=small, labelfont=bf, labelsep=period, justification=centering}
\captionsetup[sub]{font=small, labelfont=bf, labelsep=space}

% Цветовая палитра ГУАП
\usepackage{xcolor}
\definecolor{guapblue}{RGB}{0, 44, 95}       % Фирменный темно-синий ГУАП
\definecolor{guapaccent}{RGB}{0, 102, 204}   % Небесно-голубой акцент
\definecolor{darkslate}{RGB}{36, 40, 52}     % Основной текст
\definecolor{examplegreen}{RGB}{24, 115, 60} % Зеленый для примеров задач
\definecolor{amberdark}{RGB}{185, 110, 0}    % Янтарный для замечаний
\definecolor{bordergray}{RGB}{210, 215, 225}

% Выделенные блоки tcolorbox
\usepackage{tcolorbox}
\tcbuselibrary{skins,breakable}

% 1. Определение
\newtcolorbox{defbox}[1]{
    enhanced,
    breakable,
    colback=guapblue!4!white,
    colframe=guapblue,
    coltitle=white,
    fonttitle=\bfseries\sffamily,
    title=Определение: #1,
    arc=2mm,
    boxrule=1pt,
    top=3mm, bottom=3mm, left=3mm, right=3mm
}

% 2. Закон / Правило
\newtcolorbox{lawbox}[1]{
    enhanced,
    breakable,
    colback=guapaccent!5!white,
    colframe=guapaccent!85!black,
    coltitle=white,
    fonttitle=\bfseries\sffamily,
    title=Закон / Правило: #1,
    arc=2mm,
    boxrule=1pt,
    top=3mm, bottom=3mm, left=3mm, right=3mm
}

% 3. Пример решения задачи
\newtcolorbox{exbox}[1]{
    enhanced,
    breakable,
    colback=examplegreen!4!white,
    colframe=examplegreen,
    coltitle=white,
    fonttitle=\bfseries\sffamily,
    title=Пример решения задачи: #1,
    arc=2mm,
    boxrule=1pt,
    top=3mm, bottom=3mm, left=3mm, right=3mm
}

% 4. Методическое замечание / Важное указание
\newtcolorbox{notebox}[1]{
    enhanced,
    breakable,
    colback=amberdark!5!white,
    colframe=amberdark,
    coltitle=white,
    fonttitle=\bfseries\sffamily,
    title=Методическое замечание: #1,
    arc=2mm,
    boxrule=1pt,
    top=3mm, bottom=3mm, left=3mm, right=3mm
}

% 5. Контрольные вопросы и задания
\newtcolorbox{questbox}[1]{
    enhanced,
    breakable,
    colback=guapblue!3!white,
    colframe=guapblue!70!black,
    coltitle=white,
    fonttitle=\bfseries\sffamily,
    title=Задание для самопроверки: #1,
    arc=2mm,
    boxrule=0.8pt,
    top=3mm, bottom=3mm, left=3mm, right=3mm
}

% Оформление колонтитулов
\usepackage{fancyhdr}
\pagestyle{fancy}
\fancyhf{}
\fancyhead[L]{\small\sffamily\textcolor{guapblue}{Физические основы нанотехнологий}}
\fancyhead[R]{\small\sffamily\textcolor{guapblue}{Глава 2. Термодинамика наноструктур}}
\fancyfoot[C]{\sffamily\thepage}
\renewcommand{\headrulewidth}{0.5pt}
\renewcommand{\footrulewidth}{0pt}

% Ссылки и навигация
\usepackage[unicode, colorlinks=true, linkcolor=guapblue, citecolor=examplegreen, urlcolor=guapaccent]{hyperref}

% Заголовки разделов
\usepackage{titlesec}
\titleformat{\section}
  {\normalfont\Large\bfseries\sffamily\color{guapblue}}{\thesection}{1em}{}
\titleformat{\subsection}
  {\normalfont\large\bfseries\sffamily\color{guapblue!85!black}}{\thesubsection}{1em}{}
\titleformat{\subsubsection}
  {\normalfont\normalsize\bfseries\sffamily\color{guapblue!70!black}}{\thesubsubsection}{1em}{}

% Путь к рисункам
\graphicspath{{figures/}}
